
\documentclass[a4paper,12pt]{article}
\usepackage[italian]{babel}
\usepackage{listings}
\usepackage{geometry}
\geometry{margin=2.5cm}

\title{Soluzioni - Verifica SQL Bike Sharing}
\date{}

\begin{document}
\maketitle

\section*{Soluzioni Query}

\begin{enumerate}
	
	\item
	\begin{verbatim}
		SELECT nome, cognome
		FROM CLIENTE
		ORDER BY cognome, nome;
	\end{verbatim}
	
	\item
	\begin{verbatim}
		SELECT idBici, modello, tipo
		FROM BICI
		WHERE tipo = 'elettrica'
		ORDER BY modello, idBici DESC;
	\end{verbatim}
	
	\item
	\begin{verbatim}
		SELECT nome, numeroPosti
		FROM STAZIONE
		WHERE citta = 'Milano'
		ORDER BY numeroPosti DESC;
	\end{verbatim}
	
	\item
	\begin{verbatim}
		SELECT B.modello, B.tipo, S.nome
		FROM BICI B
		JOIN STAZIONE S
		ON B.idStazione = S.idStazione;
	\end{verbatim}
	
	\item
	\begin{verbatim}
		SELECT idStazione, COUNT(*) AS numero_bici
		FROM BICI
		GROUP BY idStazione;
	\end{verbatim}
	
	\item
	\begin{verbatim}
		SELECT N.idNoleggio, N.dataInizio, C.nome, C.cognome
		FROM NOLEGGIO N
		JOIN CLIENTE C
		ON N.idCliente = C.idCliente;
	\end{verbatim}
	
	\item
	\begin{verbatim}
		SELECT B.modello, C.nome, N.dataInizio
		FROM NOLEGGIO N
		JOIN CLIENTE C ON N.idCliente = C.idCliente
		JOIN BICI B ON N.idBici = B.idBici;
	\end{verbatim}
	
	\item
	\begin{verbatim}
		SELECT C.nome, C.cognome, COUNT(*) AS numero_noleggi
		FROM CLIENTE C
		JOIN NOLEGGIO N ON C.idCliente = N.idCliente
		GROUP BY C.idCliente, C.nome, C.cognome;
	\end{verbatim}
	
	\item
	\begin{verbatim}
		SELECT C.nome, C.cognome
		FROM CLIENTE C
		LEFT JOIN NOLEGGIO N
		ON C.idCliente = N.idCliente
		WHERE N.idNoleggio IS NULL;
	\end{verbatim}
	
	\item
	\begin{verbatim}
		SELECT B.idBici, B.modello
		FROM BICI B
		LEFT JOIN NOLEGGIO N
		ON B.idBici = N.idBici
		WHERE N.idNoleggio IS NULL;
	\end{verbatim}
	
	\item
	\begin{verbatim}
		SELECT idBici, COUNT(*) AS numero_noleggi
		FROM NOLEGGIO
		GROUP BY idBici;
	\end{verbatim}
	
	\item
	\begin{verbatim}
		SELECT S.nome, COUNT(B.idBici) AS numero_bici
		FROM STAZIONE S
		JOIN BICI B
		ON S.idStazione = B.idStazione
		GROUP BY S.idStazione, S.nome
		HAVING COUNT(B.idBici) > 3;
	\end{verbatim}
	
	\item
	\begin{verbatim}
		SELECT C.nome, C.cognome, COUNT(*) AS numero_noleggi
		FROM CLIENTE C
		JOIN NOLEGGIO N
		ON C.idCliente = N.idCliente
		GROUP BY C.idCliente, C.nome, C.cognome
		HAVING COUNT(*) >= 2;
	\end{verbatim}
	
	\item
	\begin{verbatim}
		SELECT DISTINCT B.modello
		FROM BICI B
		JOIN NOLEGGIO N ON B.idBici = N.idBici
		LEFT JOIN MANUTENZIONE M ON B.idBici = M.idBici
		WHERE M.idManutenzione IS NULL;
	\end{verbatim}
	
	\item
	\begin{verbatim}
		SELECT idBici, COUNT(*) AS numero_noleggi
		FROM NOLEGGIO
		GROUP BY idBici
		ORDER BY numero_noleggi DESC
		LIMIT 1;
	\end{verbatim}
	
\end{enumerate}

\section*{DDL}

\begin{verbatim}
CREATE TABLE CLIENTE (
 idCliente INT PRIMARY KEY,
 nome VARCHAR(50),
 cognome VARCHAR(50),
 email VARCHAR(100),
 citta VARCHAR(50)
);

CREATE TABLE NOLEGGIO (
 idNoleggio INT PRIMARY KEY,
 dataInizio DATE,
 dataFine DATE,
 idCliente INT,
 FOREIGN KEY (idCliente)
 REFERENCES CLIENTE(idCliente)
);
\end{verbatim}

\section*{Vista}

\begin{verbatim}
CREATE VIEW vista_noleggi_clienti AS
SELECT C.nome,
       C.cognome,
       B.modello,
       N.dataInizio
FROM NOLEGGIO N
JOIN CLIENTE C ON N.idCliente=C.idCliente
JOIN BICI B ON N.idBici=B.idBici;
\end{verbatim}

\end{document}
